\scanpage{86}
$B'_{T'}$ est obtenu à partir de $B'_T$ en regardant les flèches $H$ -- d'où catégorie, pas néc. $\in\Delta$ -- $\deleted{\unsure{B^{0}_{T'}}}\unsure{\mathcal G|G}(B'_T,H,s,b)$ (enveloppe d'un \deleted{\unsure{catég}}\inserted{type de} diagramme) \deleted{et}, en prenant la $\Delta$'enveloppe \inserted{$B^{\prime0}_{T'}$}, et en y ``rendant inv.'' \inserted{(au sens $\Delta$)} les flèches qui expriment que le composé
\[
 B_T\xrightarrow[\substack{\text{foncteur}\\\text{can (pas}\\\text{néc. un}\\\Delta\text{-foncteur)}}]{}B^{\prime0}_{T'}
 \xrightarrow[\Delta\text{-foncteur}]{}C
\]
est un $\Delta$-foncteur.

\textcircled{\unsure{$*$}} \uline{NB} Si $T$ repr. par \deleted{$R$, $T'$ repr. par $R'$ déduit de $R=\unsure{d_R|H_R}$}\inserted{et $B'_T\simeq R$, alors $T'$ est repr par $B'$} $(=H,s_R,b_R$, où $s_R,b_R$ déduits de $s,b:H\mathbin{\unsure{-|\to}}\mathrm{Ob}\,B'_{T'}$ en comp. avec

\uline{Re NB}\quad $T'\mathbin{-}T$ fidèle, donc si on a un foncteur fidèle $T\mathbin{-}\tau^I$, le composé aussi.

\textcircled{d} \inserted{composition de structures \unsure{alg.}} Si on a des \deleted{\unsure{cat}} théories $T_\alpha$, avec $B'_{T_\alpha}$, alors la théorie $T'=\prod T_\alpha:C\mapsto\prod T_\alpha(C)$ est une structure alg.
\[
 B'_{T'}=\text{2-somme amalgamée des } B'_{T_\alpha}\text{ dans }\Delta .
\]
% slanted
\margin{\uline{petit} ens. d'indices $\alpha$}

\uline{NB} Si les $T_\alpha$ repr par les $B'_{T_\alpha}$, alors $T'$ repr. par \deleted{\unsure{le}} $B'_{T'}$.

\uline{Re NB}\quad Si $T_\alpha\mathbin{-}\tau^{I_\alpha}$ fidèl, alors $T'\mathbin{-}\tau^{I'}$ fidèl, où $I=\coprod I_\alpha$
% slanted
\margin{\uline{insuffisant}}

Si on préfère ne pas changer en cours de route des objets de base, on posera

\textcircled{$d'$} \inserted{(Composition de structures avec \deleted{\unsure{objet|objets}} bases fixées)} Si on a des théories $T_\alpha$ \deleted{[fidèles sur $\tau^I$]} $\varphi_\alpha:T_\alpha\mathbin{-}\tau^I$, alors le produit 2-fibré des $T_\alpha$ sur $\tau^I$
\[
 T'(C)=\prod_{\alpha\unsure{\,/\,|\,,\,}C^I}^{(2)}T_\alpha(C)
 \qquad\deleted{\text{(fibrés sur }\unsure{-}}
\]
\deleted{\unsure{est}} On pose $B'_{T'}=$ \deleted{\unsure{caté}} somme amalgamée des $B'_{T_\alpha}$ sur $B'_{\tau^I}=\widetilde I$.

\uline{NB} Si les $T_\alpha$ repr par les $B'_{T_\alpha}$, $T$ l'est par $B_T$.

\uline{Re NB} \deleted{\unsure{Si les}} On peut vouloir se restreindre \unsure{aux|au x} $T_\alpha\mathbin{-}\tau^{\unsure{I_\alpha|I}}$ fidèles, alors $T'\to\tau^I$ est fidèle.
% archive sheet number: 85
